
\begin{frame}{Appendix --- Consolidation Robust to the Full Pre-Treatment Path}
\hypertarget{app:ptrobust}{}
\small
The balance test clears the margin, effective-N and top-two share, with the Herfindahl borderline. As a further check we add the 2020 base-year level to the headline 2016 lag, so the two lags $\{2016, 2020\}$ span both where a municipality started and the pre-slope it was on. The 2020 level is pre-determined with respect to the 2020$\to$2024 shock, since the instrument's shares are dated 2020 and the shifter leaves the own state out, so it is a legitimate control. The coefficient is then identified off the 2024 deviation from each municipality's own path.
\smallskip
\begin{center}
\scriptsize
\resizebox{\linewidth}{!}{% Auto-generated by code/03_estimation/02_iv_main.R
% Do not edit manually -- rerun the generating script to update

\begin{tabular}{l*{4}{c}}
\toprule\toprule
Dep.\ var.: & \textbf{$\Delta$ Margin (W$-$RU)} & \textbf{$\Delta$ Eff.\ N candidates} & \textbf{$\Delta$ Vote HHI} & \textbf{$\Delta$ Top-2 vote share} \\
\midrule
\rowcolor{mylight}
\textbf{Judicialization} & 0.057$^{**}$ & -0.044 & 0.022 & 0.001 \\
\rowcolor{mylight}
 & \textcolor{mygray}{(0.025)} & \textcolor{mygray}{(0.040)} & \textcolor{mygray}{(0.013)} & \textcolor{mygray}{(0.009)} \\
\midrule
$N$ & 5,560 & 5,560 & 5,560 & 5,560 \\
2024 Mean & 0.270 & 2.021 & 0.527 & 0.945 \\
2016 Mean & 0.193 & 2.172 & 0.492 & 0.929 \\
\bottomrule\bottomrule
\end{tabular}
}
\end{center}
\smallskip
\footnotesize The estimates are essentially unchanged. The top-two margin moves from its headline value $\MarginCoef$ to $\PtRobMarginCoef$ ($p=\PtRobMarginP$), while effective-N and the Herfindahl shift only in the third decimal and remain imprecise, as in the headline. Conditioning on both the 2020 level and the pre-slope a municipality was on leaves the result intact, which places the consolidation as a response to the 2020$\to$2024 shock rather than the continuation of a pre-existing path. \hfill\hyperlink{src:pretrend}{\beamerreturnbutton{Back to balance test}}
\end{frame}
